\subsection{商律}

定义11. 设E为一个集合。一个合成律 \(\top\) 与E上的一个等价关系R称为相容的，如果关系 \(x \equiv x'\) （模R）和 \(y \equiv y'\) （模R）（对于 \(x, x', y, y' \in \mathbf{E}\) ）蕴含 \(x \top y \equiv x' \top y'\) （模R）；商集E/R上的合成律将 \(x\) 和 \(y\) 的等价类映至 \(x \top y\) 的等价类，称为律 \(\top\) 关于R的商律。带有商律的集合E/R称为E关于R的商广群。

要说明 E 上的等价关系 R 与 E 上的内部合成律 \(f: E \times E \to E\) 相容，指的是映射 f 与乘积等价关系 \(R \times R\) （在 \(E \times E\) 上）以及等价关系 R（在 E 上）相容（在集合论，II，§6，第5号的意义下）。（集合论，II，§6，第8号）。这也就意味着 R 的图是 \(E \times E\) 的一个子广群.

如果律 \(\top\) 是结合的（相应地，交换的），那么商律也是结合的（相应地，交换的）（更简洁地说，我们称结合性或交换性在过渡到商时得以保持）。

从广群 E 到商广群 E/R 的典范映射是一个同态。

对于从 E/R 到广群 F 的映射 g，要使 g 成为同态，其充要条件是 g 与从 E 到 E/R 的典范映射的复合是一个同态。

以下两个命题由定义直接可得：

命题6. 设 E 和 F 是两个广群，f 是 E 到 F 的一个同态。设 R\{x,y\} 表示关系 \(f(x)=f(y)\) （在 E 的元素 x, y 之间）。则 R 是 E 上的一个等价关系，与 E 上的律相容，并且由 f 通过取商导出的从 E/R 到 f(E) 的映射是商广群 E/R 到 F 的子广群 f(E) 上的一个同构。

命题7. 设 E 是一个广群，R 是 E 上与 E 上的律相容的一个等价关系。对于 E/R 上的一个等价关系 S，它与商律相容的充要条件是 S 具有 T/R 的形式，其中 T 是 E 上的一个由 R 蕴含且与 E 上的律相容的等价关系。此时，从 E/T 到 (E/R)/(T/R) 的典范映射（集合论，II，§6，第7号）是一个广群同构。

命题8. 设E是一个广群，A是E的一个稳定子集，R是E上与E的律相容的一个等价关系。A关于R的饱和B（集合

理论，II，§6，第5号）是一个稳定子集。等价关系 \(\mathbf{R}_{\mathbf{A}}\) 和 \(\mathbf{R}_{\mathbf{B}}\) （分别由 \(\mathbf{R}\) 在 \(\mathbf{A}\) 和 \(\mathbf{B}\) 上诱导）与诱导律相容，且由典范嵌入 \(\mathbf{A}\) 到 \(\mathbf{B}\) 通过取商得到的映射是 \(\mathbf{A} / \mathbf{R}_{\mathbf{A}}\) 到 \(\mathbf{B} / \mathbf{R}_{\mathbf{B}}\) 的一个广群同构.

设 \(\top\) 表示 E 上的律。如果 \(x\) 和 \(y\) 是 B 的两个元素，则存在两个元素 \(x'\) 和 \(y'\) （属于 A ）使得 \(x \equiv x' \pmod{R}\) 和 \(y \equiv y' \pmod{R}\) ；于是 \(x \top y \equiv x' \top y' \pmod{R}\) 且 \(x' \top y' \in A\) ，从而 \(x \top y \in B\) 。因此，B 是 E 的一个稳定子集，其余断言是显然的。

设 M 是一个广群，且 \(((u_{\alpha}, v_{\alpha}))_{\alpha \in I}\) 为 \(M \times M\) 的元素的一个族。考虑 M 上所有与 M 上的律相容且满足 \(u_{\alpha} \equiv v_{\alpha} \pmod{S}\) （对所有 \(\alpha \in I\) ）的等价关系 S。这些关系的图的交集是一个等价关系 R 的图，该关系与 M 上的律相容且满足 \(u_{\alpha} \equiv v_{\alpha} \pmod{R}\) 。因此，R 是具有这两个性质的最细（集合论，III，§1，第3和7节）等价关系。它称为与 M 上的律相容且由 \((u_{\alpha}, v_{\alpha})\) 所生成的等价关系.

命题9. 保持上述记号，设 \(f\) 是 \(\mathbf{M}\) 到一个广群的同态，使得 \(f(u_{\alpha}) = f(v_{\alpha})\) 对所有 \(\alpha \in \mathbf{I}\) 成立。则 \(f\) 与 \(\mathbf{R}\) 相容.

设 T 为与 f 相关联的等价关系。则 \(u_{\alpha} \equiv v_{\alpha} \pmod{T}\) 对所有 \(\alpha \in I\) 成立，且 T 与 M 上的律相容，因此 T 比 R 更粗；这就证明了该命题。

